\begin{pastexam}[label={px:ncku-114-isec-ds-11}]{成大智慧資安學程 114 學年度 單選題第 11 題}
（單選；每題 5 分，答錯倒扣 2 分）

Given a hash table with 11 buckets; labeled from 0 to 10. The hash function h(key)= key \% 11 allocates
data into a hash table. The chaining strategy is used to store data in each bucket. Please show the total
number of comparisons required to store the numbers ``3, 41, 15, 36, 74, 58, 91, 45, 48, 64''.

a)~20 \qquad b)~16 \qquad c)~22 \qquad d)~14 \qquad e)~None of the above.
\end{pastexam}

\begin{pastsoln}
答案是 e) None of the above。

\textbf{題目的歧義。} 「存入時的比較次數」沒有定義怎麼算。 鏈結法插入時，通常要先在那條串列裡找一遍，
確認沒有重複（或找到串列尾端再接上），所以比較次數就是插入當時串列裡已經有的元素數（資結書
第~12.2~節）。 也有人把「檢查到串列尾」或「計算雜湊後看桶子」各算一次。 下面先算雜湊值，再用兩種
算法算。

\begin{center}
\small
\begin{tabular}{rcc}
\toprule
key & key \% 11 & 放進去之前，桶子裡已有幾個 \\
\midrule
3 & 3 & 0 \\
41 & 8 & 0 \\
15 & 4 & 0 \\
36 & 3 & 1 \\
74 & 8 & 1 \\
58 & 3 & 2 \\
91 & 3 & 3 \\
45 & 1 & 0 \\
48 & 4 & 1 \\
64 & 9 & 0 \\
\bottomrule
\end{tabular}
\end{center}
最後桶子 $3$：$3, 36, 58, 91$；桶子 $8$：$41, 74$；桶子 $4$：$15, 48$；桶子 $1$：$45$；桶子 $9$：$64$。
\begin{itemize}
\item 只算和已存在元素的比較：$1 + 1 + 2 + 3 + 1 = 8$ 次。
\item 每次插入再多算一次（例如和空串列或串列尾的那一次）：$8 + 10 = 18$ 次。
\end{itemize}
兩種常見算法得到 $8$ 或 $18$，都不在 a)--d) 裡，答案 e)。 （若把成功搜尋這十個數的比較次數加起來：
桶子 $3$ 是 $1 + 2 + 3 + 4$，桶子 $8$、$4$ 各是 $1 + 2$，其餘各 $1$，共 $10 + 3 + 3 + 2 = 18$，也是 $18$。）
\end{pastsoln}
